% Source: https://artofproblemsolving.com/community/c3h1721671_1995_aime_5
% Stable identity: aops-topic-1721671-q01
Because the coefficients are real, the non-real roots come in complex-conjugate pairs: the conjugate of every root is also a root (with the same multiplicity).

Let $z,w$ be the two roots with $zw=13+i$. They are not conjugates of each other, since $z\bar z=|z|^2$ is real and $13+i$ is not. So the conjugates $\bar z$ and $\bar w$ are the other two roots. This is why their product is exactly $\bar z\,\bar w=\overline{zw}=13-i$ (not some multiple of it): the other two roots are \emph{the conjugates of} $z$ and $w$. Likewise $\bar z+\bar w=3+4i$, so $z+w=3-4i$.

By Vieta's formulas, $b$ is the sum of the products of the roots taken two at a time:
\[
b=zw+\bar z\bar w+(z+w)(\bar z+\bar w)=(13+i)+(13-i)+(3-4i)(3+4i)=26+25=\boxed{51}.
\]
